%!TEX TS-program = xelatex
%!TEX encoding = UTF-8 Unicode
\documentclass[11pt,a4paper]{awesome-cv}
\usepackage[hidelinks]{hyperref}

% --- ATS / PDF metadata ---
\hypersetup{
  pdftitle={Juan Manuel Young Hoyos — Cover Letter — Software Engineer, Frontend},
  pdfauthor={Juan Manuel Young Hoyos},
  pdfsubject={Cover Letter — Software Engineer, Frontend — Helsing},
  pdfkeywords={React, TypeScript, Rust, component library, design system, frontend infrastructure, developer tooling, data visualization, Rerun, frontend standards, Vite, REST, GraphQL, Rust, performance, Helsing}
}

\geometry{left=1.4cm, top=.8cm, right=1.4cm, bottom=1.8cm, footskip=.5cm}

% Emit ActualText so PDF text extraction yields ASCII hyphens (ATS keyword matching)
\XeTeXgenerateactualtext=1

\colorlet{awesome}{awesome-red}
\setbool{acvSectionColorHighlight}{true}
\renewcommand{\acvHeaderSocialSep}{\quad\textbar\quad}

\definecolor{darktext}{HTML}{1A1A1A}
\colorlet{text}{darktext}
\renewcommand*{\lettertextstyle}{\fontsize{10pt}{1.45em}\bodyfontlight\upshape\color{darktext}}
\renewcommand*{\lettersectionstyle}[1]{{\fontsize{14pt}{1em}\bodyfont\bfseries\color{darktext}#1}}

% --- Header info ---
\name{Juan Manuel}{Young Hoyos}
\position{Software Engineer, Frontend — React \& TypeScript · Component Libraries \& Standards}
\address{Medellín, Colombia · Open to relocating to Munich / Berlin / London · English B2+}
\mobile{+57 1234567890}
\email{juanmanuel12.13jmyh81@gmail.com}
\homepage{jmyounghoyos.com}
\github{Youngermaster}
\linkedin{juan-manuel-young-hoyos}

% --- Letter meta ---
\recipient
  {Frontend Engineering — Hiring Team}
  {Helsing}
\letterdate{\today}
\lettertitle{Application for Software Engineer — Frontend}
\letteropening{Dear Helsing Hiring Team,}
\letterclosing{Sincerely,}
\letterenclosure[Attached]{Resume — Juan Manuel Young Hoyos.pdf}

\begin{document}
\makecvheader[R]
\makelettertitle

\begin{cvletter}

I am applying for the \textbf{Frontend Software Engineer} role at Helsing. What stands out is that you have relatively few frontend engineers, which means the work is foundational: shaping the \textbf{React component library}, the standards and the tooling that everyone else builds on. That kind of leverage is what I look for, and the seriousness of the domain --- including the ethical debate you invite --- is a reason I want to work on it rather than a reservation.

I write \textbf{React} and \textbf{TypeScript} for live production software. As Technical Lead at \textbf{Okorum Technologies} I own the front end and the services behind it for a national-scale platform delivered ahead of a hard deadline, structured as a \textbf{Turborepo} monorepo where shared types and one CI pipeline keep contracts honest across front ends --- the same instinct that makes shared component foundations worth building.

I have measurable results on the things you care about. I cut a React/Angular bundle from \textbf{3.46 MB to ~886 KB}, built \textbf{reusable TypeScript components} that lifted Wegmans in-app engagement \textbf{30\%}, cut critical bugs \textbf{25\%} with a Jest testing strategy and Dynatrace monitoring, and led an e-commerce redesign for Duluth Trading Company in \textbf{Next.js} with responsive, accessible UX. I have also raised TypeScript quality across other teams' React Native codebases --- unglamorous work that makes everyone faster.

Your posting says Rust is not a prerequisite, but it happens to be where my own time goes. I run \textbf{Project Olivaw} (github.com/Project-Olivaw), robotics libraries in pure \textbf{Rust}: \textbf{olivaw-slam}, a 2D lidar SLAM implementation; \textbf{olivaw-lidar}, an RPLIDAR driver with no C++ SDK or FFI; and \textbf{olivaw-cli}, a developer tool that vendors editable components into a project from a curated registry --- a component library engineers own rather than depend on, which is the same instinct I would bring to your React component library. I have also shipped Go, Python, Node.js and C++ (including a WebAssembly codebase) and work with \textbf{REST and GraphQL} contracts daily, so navigating an unfamiliar backend to unblock myself is routine.

The day-to-day you describe suits me: reviewing colleagues' work before advancing my own, writing RFCs, collaborating asynchronously, and mentoring. I have led engineering and QA teams and taught for over a decade. Thank you for your consideration.

\end{cvletter}

\makeletterclosing
\end{document}
